\subsection{向量空间的张量积}

第3、4、5节中关于自由模或投射模的结果特别适用于向量空间，并给出以下性质：

命题14. 给定一个正合序列

\[
0 \rightarrow \mathbf {E} ^ {\prime} \rightarrow \mathbf {E} \rightarrow \mathbf {E} ^ {\prime \prime} \rightarrow 0\tag{23}
\]

关于域 K 上的右向量空间与线性映射，以及域 K 上的左向量空间 F，相应的 Z-线性映射序列

\[
0 \rightarrow \mathrm{E} ^ {\prime} \otimes_ {\mathbb {K}} \mathrm{F} \rightarrow \mathrm{E} \otimes_ {\mathbb {K}} \mathrm{F} \rightarrow \mathrm{E} ^ {\prime \prime} \otimes_ {\mathbb {K}} \mathrm{F} \rightarrow 0
\]

是正合的且分裂。

由于序列 (23) 分裂，这是 §3 第 7 号命题 7 的推论 5 以及 §3 第 6 号命题 5 的一个特例。

由于命题 14，当 \(\mathbf{E}'\) 是 \(\mathbf{E}\) 的一个向量子空间， \(j: \mathbf{E}' \to \mathbf{E}\) 为典范单射时， \(\mathbf{E}' \otimes_{\mathbb{K}} \mathbf{F}\) 通常作为 \(\mathbf{Z}\) -模被等同于 \(\mathbf{E} \otimes_{\mathbb{K}} \mathbf{F}\) 的一个子模，所借助的是单射 \(j \otimes 1_{\mathbf{F}}\) 。在此约定下：

推论。设 K 为域，E 为 K 上的右向量空间，F 为 K 上的左向量空间， \((\mathbf{M}_{\alpha})_{\alpha \in \mathbf{A}}\) 为 E 的一族向量子空间， \((\mathbf{N}_{\beta})_{\beta \in \mathbf{B}}\) 为 F 的一族向量子空间。则

\[
\left(\bigcap_ {\alpha \in \mathrm{A}} \mathrm{M} _ {\alpha}\right) \otimes_ {\mathrm{K}} \left(\bigcap_ {\beta \in \mathrm{B}} \mathrm{N} _ {\beta}\right) = \bigcap_ {(\alpha , \beta) \in \mathrm{A} \times \mathrm{B}} \left(\mathrm{M} _ {\alpha} \otimes_ {\mathrm{K}} \mathrm{N} _ {\beta}\right).\tag{24}
\]

显然只需证明该特例

\[
\left(\bigcap_ {\alpha \in \mathrm{A}} \mathrm{M} _ {\alpha}\right) \otimes_ {\mathrm{K}} \mathrm{F} = \bigcap_ {\alpha \in \mathrm{A}} \left(\mathrm{M} _ {\alpha} \otimes_ {\mathrm{K}} \mathrm{F}\right).\tag{25}
\]

显然，(25) 的左端包含于右端。为证明反向包含关系，我们考虑 F 的一组基 \((f_{\lambda})_{\lambda \in L}\) 。 \(\mathbf{E} \otimes_{\mathbb{K}} \mathbf{F}\) 的每个元素于是都可以唯一地表示为 \(\sum_{\lambda \in L} x_{\lambda} \otimes f_{\lambda}\) 的形式，其中 \(x_{\lambda} \in \mathbf{E}\) （§ 3，第 7 号，命题 7 的推论 1）；如果 \(\mathbf{E}'\) 是 E 的一个向量子空间，则关系式 \(\sum_{\lambda \in L} x_{\lambda} \otimes f_{\lambda} \in \mathbf{E}' \otimes_{\mathbb{K}} \mathbf{F}\) 依命题 14 等价于 \(x_{\lambda} \in \mathbf{E}'\) 对所有 \(\lambda \in L\) 成立。说 \(\sum_{\lambda \in L} x_{\lambda} \otimes f_{\lambda}\) 属于每一个 \(\mathbf{M}_{\alpha} \otimes_{\mathbb{K}} \mathbf{F}\) ，这就意味着对所有 \(\lambda \in L\) 以及所有 \(\alpha \in A\) ， \(x_{\lambda} \in M_{\alpha}\) ，即 \(x_{\lambda} \in \bigcap_{\alpha \in A} M_{\alpha}\) 对所有 \(\lambda \in L\) 成立，这证明了 (25) 的右端包含于左端。命题 15. 若 \((E_{\lambda})_{\lambda \in L}\) 是域 K 上的一族右向量空间，且 \((F_{\mu})_{\mu \in M}\) 是 K 上的一族左向量空间，则典范映射

\[
\left(\prod_ {\lambda \in L} E _ {\lambda}\right) \otimes_ {K} \left(\prod_ {\mu \in M} F _ {\mu}\right)\rightarrow \prod_ {(\lambda , \mu) \in L \times M} (E _ {\lambda} \otimes_ {K} F _ {\mu})\tag{26}
\]

(§ 3，第 7 号，公式 (22)) 是单射的。

我们写 \(\mathbf{F} = \prod_{\mu \in \mathbf{M}} \mathbf{F}_{\mu}\) ；映射 (26) 是诸典范映射的复合。

\[
\left(\prod_ {\lambda \in L} \mathbf {E} _ {\lambda}\right) \otimes_ {\mathbb {K}} F \rightarrow \prod_ {\lambda \in L} \left(\mathbf {E} _ {\lambda} \otimes_ {\mathbb {K}} F\right)
\]

和

\[
\prod_ {\lambda \in L} \left(\mathrm{E} _ {\lambda} \otimes_ {K} F\right)\rightarrow \prod_ {\lambda \in L} \left(\prod_ {\mu \in M} \left(\mathrm{E} _ {\lambda} \otimes_ {K} F _ {\mu}\right)\right);
\]

由于 F 和诸 \(\mathbf{E}_{\lambda}\) 是 K 上的向量空间，这就归结为 § 3，第 7 号，命题 7 的推论 3。

当命题 15 的条件满足时，张量积 \(\left(\prod_{\lambda \in L} E_{\lambda}\right) \otimes_{K} \left(\prod_{\mu \in M} F_{\mu}\right)\) 通常等同于它在以下空间中的典范像

\[
\prod_ {\lambda , \mu} (\mathbf {E} _ {\lambda} \otimes_ {K} \mathbf {F} _ {\mu}).
\]

按照这个约定：

推论。设 F 是 K 上的左向量空间；对于每个集合 X，左向量空间 \(K_{d}^{X} \otimes_{K} F\) 等同于空间 \(F^{X}\) （即由 X 到 F 的所有映射组成的空间）的子空间，它由满足以下条件的映射 u 组成： \(u(X)\) 在 F 中具有有限秩。

若 \((f_{\lambda})\) 是 \(F\) 的一组基，则 \(\sum_{\lambda \in L} v_{\lambda} \otimes f_{\lambda}\) 作为 \(K_d^X \otimes_K F\) 的元素在 (26) 下等同于映射 \(x \mapsto \sum_{\lambda} v_{\lambda}(x)f_{\lambda}\) 。由于 \(v_{\lambda} = 0\) ，除了指标 \(\lambda\) 属于一个有限子集 \(H\) （该子集含于 \(L\) ）以外， \(X\) 在上述映射下的像包含于 \(F\) 中由诸 \(f_{\lambda}\) （其指标 \(\lambda \in H\) ）生成的有限维子空间。反之，设 \(u: X \to F\) 为一个映射，使得 \(u(X)\) 包含于一个有限维子空间 \(G\) （该子空间含于 \(F\) ），并设 \((b_i)_{1 \leqslant i \leqslant n}\) 是 \(G\) 的一组基。对所有 \(x \in X\) ，我们可以写成 \(u(x) = \sum_{i=1}^{n} v_i(x)b_i\) ，其中诸 \(v_i(x)\) 是 \(K\) 的完全确定的元素；因此 \(n\) 个映射 \(v_i: X \to K\) 被定义，且显然 \(u\) 随后与元素 \(\sum_{i=1}^{n} v_i \otimes b_i\) 等同.

类似地，对于域 K 上的右向量空间 E 和集合 Y， \(E \otimes_{K} K_{s}^{Y}\) 等同于空间 \(E^{Y}\) 的一个子空间，它由满足以下条件的映射 \(v: Y \to E\) 组成： \(v(Y)\) 具有有限秩。更特别地，对于每个域 K， \(K_{d}^{X} \otimes_{K} K_{s}^{Y}\) 等同于空间 \(K^{X \times Y}\) （由 \(X \times Y\) 到 K 中的映射组成）的一个子空间（K 被视为 (K, K)-双模）；元素 \(\sum_{i} u_{i} \otimes v_{i}\) （其中 \(u_{i}\) 是 X 到 K 的映射，而 \(v_{i}\) 是 Y 到 K 的映射）被等同于映射 \((x, y) \mapsto \sum_{i} u_{i}(x) v_{i}(y)\) ，它是 \(X \times Y\) 到 K 中的映射。

命题 16. (i) 设 K, L 为两个域，E 为 K 上的左向量空间，F 为 L 上的左向量空间，G 为 (K, L)-双模。则典范 Z-同态

\[
\nu : \operatorname{Hom} _ {\mathbf {K}} (\mathrm{E}, \mathrm{G}) \otimes_ {\mathrm{L}} \mathrm{F} \rightarrow \operatorname{Hom} _ {\mathbf {K}} (\mathrm{E}, \mathrm{G} \otimes_ {\mathrm{L}} \mathrm{F})\tag{27}
\]

（§ 4，第 2 号，公式 (7)）是单射；当向量空间 E, F 中有一个是有限维时，它是双射。

\begin{enumerate}
\def\labelenumi{(\roman{enumi})}
\setcounter{enumi}{1}
\tightlist
\item
  设 \(E_{1}\) , \(E_{2}\) , \(F_{1}\) , \(F_{2}\) 为交换域 K 上的四个向量空间；则典范 K-同态
\end{enumerate}

\[
\lambda : \operatorname{Hom} \left(\mathrm{E} _ {1}, \mathrm{F} _ {1}\right) \otimes \operatorname{Hom} \left(\mathrm{E} _ {2}, \mathrm{F} _ {2}\right)\rightarrow \operatorname{Hom} \left(\mathrm{E} _ {1} \otimes \mathrm{E} _ {2}, \mathrm{F} _ {1} \otimes \mathrm{F} _ {2}\right)\tag{28}
\]

(§ 4，第 4 号，公式 (21)) 是单射；若有序对 \((\mathbf{E}_1, \mathbf{E}_2), (\mathbf{E}_1, \mathbf{F}_1), (\mathbf{E}_2, \mathbf{F}_2)\) 之一由有限维空间组成，则它是双射。

断言 (i) 是 § 4 第 2 号命题 2 的特殊情形。类似地，(ii) 的第二个断言是 § 4 第 4 号命题 4 的特殊情形。最后，要看出同态 (28) 总是单射的，注意 \(\operatorname{Hom}(\mathbf{E}_i, \mathbf{F}_i)\) 是 \(\mathbf{F}_i^{\mathbf{E}_i}\) 的一个向量子空间 ( \(i = 1, 2\) )，并且

\[
\operatorname{Hom} \left(\mathrm{E} _ {1} \otimes \mathrm{E} _ {2}, \mathrm{F} _ {1} \otimes \mathrm{F} _ {2}\right)
\]

典范地等同于空间 \((\mathbf{F}_{1} \otimes \mathbf{F}_{2})^{\mathbf{E}_{1} \times \mathbf{E}_{2}}\) 的一个向量子空间（II，§ 3，第 1 号，命题 1）；当完成这些等同，并且 (28) 的左端也被等同于 \(F_{1}^{E_{1}} \otimes F_{2}^{E_{2}}\) 的一个子空间（命题 14）时，典范映射 (28) 就变为典范映射 (26) 在该子空间上的限制，并且已经看到（命题 15）后者是单射的。

推论 1. 设 E 和 F 是域 K 上的两个向量空间；则典范映射

\[
\mathbf {E} ^ {*} \otimes_ {\mathbb {K}} \mathbf {F} \rightarrow \operatorname{Hom} _ {\mathbb {K}} (\mathbf {E}, \mathbf {F})
\]

(§ 4，第 2 号，公式 (11)) 是单射；当 E 或 F 为有限维时，它是双射。

这是命题 16 (i) 的一个特例。

推论 2. 设 E 是右向量空间，F 是同一域 K 上的左向量空间；则典范映射

\[
\mathbf {E} \otimes_ {\mathbb {K}} \mathbf {F} \rightarrow \operatorname{Hom} _ {\mathbb {K}} (\mathbf {E} ^ {*}, \mathbf {F})\tag{29}
\]

（§ 4，第 2 号，公式 (15)）是单射；当 E 是有限维时，它是双射。

这是 § 4 第 2 号注记 2 的一个特例。

注记。(1) 设 K 是一个交换域，E、F 是 K 上的两个向量空间，

\((a_{\lambda})\) 是 \(\mathbf{E}\) 的一组基， \((b_{\mu})\) 是 \(\mathrm{F}\) 的一组基；于是 \((a_{\lambda} \otimes b_{\mu})\) 是向量 \(\mathbf{K}\) -空间 \(\mathbf{E} \otimes_{\mathbb{K}} \mathbf{F}\) 的一组基（II，§ 3，第 7 号，命题 2 的推论 2），因此

\[
\dim_ {\mathbf {K}} (\mathbf {E} \otimes_ {\mathbf {K}} \mathbf {F}) = \dim_ {\mathbf {K}} \mathbf {E}. \dim_ {\mathbf {K}} \mathbf {F}.\tag{30}
\]

\begin{enumerate}
\def\labelenumi{(\arabic{enumi})}
\setcounter{enumi}{1}
\tightlist
\item
  设 K 是一个交换域， \(E_{1}\) , \(E_{2}\) , \(F_{1}\) , \(F_{2}\) 为 K 上的四个向量空间， \(u: E_{1} \rightarrow F_{1}\) , \(v: E_{2} \rightarrow F_{2}\) 为两个线性映射；则
\end{enumerate}

\[
\operatorname{rg} (u \otimes v) = \operatorname{rg} (u). \operatorname{rg} (v).\tag{31}
\]

显然 \((u\otimes v)(\mathbf{E}_1\otimes \mathbf{E}_2)\) 是 \(u(\mathbf{E}_1)\otimes v(\mathbf{E}_2)\) 在 \(\mathbf{F}_1\otimes \mathbf{F}_2\) 中的典范像，因此（命题 14）同构于 \(u(\mathbf{E}_1)\otimes v(\mathbf{E}_2)\) ；结论随后由 (30) 得出。

\begin{enumerate}
\def\labelenumi{(\arabic{enumi})}
\setcounter{enumi}{2}
\tightlist
\item
  在注记 1 的相同假设下，
\end{enumerate}

\[
\dim_ {\mathbb {K}} (\operatorname{Hom} _ {\mathbb {K}} (\mathrm{E}, \mathrm{F})) \geqslant \dim_ {\mathbb {K}} \mathrm{E}. \dim_ {\mathbb {K}} \mathrm{F}.\tag{32}
\]

如果 E 同构于 \(K^{(I)}\) ，则 \(\operatorname{Hom}(E, F)\) 同构于 \((\operatorname{Hom}(K, F))^{I}\) （§ 1，第 6 号，命题 6 的推论 1），从而同构于 \(F^{I}\) （§ 1，第 14 号）；由于 \(F^{(I)}\) 是 \(F^{I}\) 的一个子空间，且 \(\dim(F^{(I)}) = \operatorname{Card}(I).dim F = \dim E. dim F\) （第 2 号，命题 2），不等式 (32) 由第 3 号命题 4 的推论 4 得出。同样的论证表明，当 dim E 有限时，(32) 的两边相等（参见 § 10，第 3 和第 4 号）。

